\section{Parametry procesu uczenia}

Poza architekturą i współczynnikiem uczenia pozostałe elementy procesu
uczenia są utrzymywane na stałym poziomie, aby porównanie konfiguracji było
miarodajne. Do takich elementów należą:
\begin{itemize}
    \item rozmiar mini-batcha,
    \item liczba epok,
    \item zastosowana funkcja kosztu,
    \item zastosowany optymalizator,
    \item sposób podziału danych i standaryzacji cech.
\end{itemize}

\subsection{Stałe parametry eksperymentów}

Podczas porównywania konfiguracji wykorzystano następujące ustawienia:
\begin{itemize}
    \item funkcja kosztu: \texttt{Cross Entropy Loss},
    \item optymalizator: \texttt{Stochastic Gradient Descent},
    \item rozmiar mini-batcha: \(16\),
    \item liczba epok: \(1000\),
    \item ewaluacja wykonywana co \(100\) epok.
\end{itemize}

\subsection{Powtarzanie eksperymentów}

Każda konfiguracja może zostać uruchomiona wielokrotnie dla różnych wartości
\texttt{random\_state}. Dzięki temu możliwe jest porównanie nie tylko
pojedynczego wyniku, ale również średnich wartości metryk oraz ich zmienności.
